% Staged author update. Existing article inputs remain unchanged until activation.
\subsection*{Reference-image history intervention}
The unchanged balanced checkpoint receives each of the 64 normal physical QA
prompts once, with no optimization or external native decoding. After prompt
prefill, the next-writing state is computed at every complete reference boundary
using the canonical preceding glyph images. The supplied reference length
determines the number of diagnostic steps; the stopping predictor is observed
without terminating the trace. These teacher-forced traces must not be combined
with autonomous answer counts. Source images, initial states and the bridge
remain unchanged.

Each layout contains 745 target glyph positions and 32 terminal boundaries.
The first position of each case is reported separately, leaving 713 later
positions. Foreground F1 is $2\mathrm{TP}/(2\mathrm{TP}+\mathrm{FP}+\mathrm{FN})$,
with value one for two empty ink masks. Exact equality compares the complete
binary bitmap. Punctuation and spaces remain in the pools. Summaries aggregate
positions unless explicitly identified as equal-case means. Shared autonomous
and reference indices are retained only where both predictions exist, with
missing and extra positions reported separately. A shared numeric index is not
semantic alignment when wording or output lengths differ.

The later-writing gate requires both median F1 at least 0.99 and at least
90\% of positions reaching 0.95. The stopping diagnosis separately requires
the prospectively fixed nonterminal and terminal rates. All 1,490 reference
positions, initial-state comparisons and saved autonomous pixel/stop replays
are independently checked. Review then inspects all 8,561 displayed tiles.
The common-prefix state comparison is post hoc and descriptive. All cases are
reused development data, with reviewer familiarity and no independent human
adjudication. The diagnosis takes 460.834 seconds, including 70.067 seconds
pressure waiting, at 17.786\,GiB peak allocated GPU memory. Loading is 137.167
seconds; loading and waiting intervals may overlap. Its 1,554 contextual calls
do not create new autonomous answers or change original language scores.

\subsection*{Immutable prompt memory and visual continuation}
Let $X$ be the physical prompt image and $Y_0,\ldots,Y_{N-1}$ the complete target
glyph images. Each signed image has shape $1\times32\times32$, with ink at
$-1$ and background at $+1$. Source prompt prefill produces detached memory
$C_\ell(X)=(K^P_\ell,V^P_\ell)$ for each native layer
$\ell=0,\ldots,35$. The new residuals are bypassed during this computation.
Source vision, contextual parameters and the original zero-output source
adapters are frozen. The answer input sequence is
\begin{equation}
 a_0=q,\qquad a_{t+1}=S(E(Y_t)),\quad 0\leq t<N,
 \label{eq:continuation-input}
\end{equation}
where $E$ is the frozen image encoder and $S$ its fixed affine denormalization
into the imported 4,096-dimensional contextual coordinates. The learned query
$q$ initially copies the final fixed numeric task-control vector. It contains
no case-specific answer or character identity.

With $h^0_t=a_t$, the answer transition is
\begin{align}
 u^\ell_t&=F_\ell(h^\ell_t;C_\ell(X),h^\ell_{<t}),\nonumber\\
 h^{\ell+1}_t&=u^\ell_t+\frac{\alpha}{r}B_\ell A_\ell u^\ell_t.
 \label{eq:continuation-residual}
\end{align}
Here $F_\ell$ is the frozen native layer,
$A_\ell\in\mathbb R^{8\times4096}$,
$B_\ell\in\mathbb R^{4096\times8}$, $r=8$ and $\alpha=16$.
Each $A_\ell$ uses Kaiming-uniform initialization and each $B_\ell$ begins at
zero. The residual is applied after the native layer, including to the query.
Prompt and answer position offsets preserve the imported rotary convention.
The causal mask allows the complete source prompt and current/preceding answer
inputs. Thus state $t$ predicts $Y_t$ before that target image enters the history.
Final native normalization produces $z_t$.

Only frozen prompt keys and values are cached during fitting. The cache update
is a pure concatenation of prompt memory with the current answer keys/values;
activation recomputation cannot append to or lengthen the prompt. All 36 answer
layers are recomputed with their current residuals using nonreentrant
checkpointing. Incremental inference maintains a separate answer cache for
the current output stream and leaves the prompt memory untouched. No mutable
answer cache is captured during gradient-enabled execution.

\subsection*{Physical writing loss and autonomous feedback}
The trainable readout $R$ and frozen decoder $D$ predict pixel logits
$L_t=D(R(z_t))$. A separate trainable head predicts stopping logit $s_t$.
Writing $b_{ti}=(1-Y_{ti})/2$ for reference ink at pixel $i$, the implemented
balanced reconstruction objective is
\begin{equation}
 \begin{split}
 \mathcal L_{\rm image}=\frac12\biggl[
 &\frac{\sum_{t,i}b_{ti}\,\operatorname{softplus}(-L_{ti})}
 {\max(1,\sum_{t,i}b_{ti})}\\
 +&\frac{\sum_{t,i}(1-b_{ti})\,\operatorname{softplus}(L_{ti})}
 {\max(1,\sum_{t,i}(1-b_{ti}))}\biggr].
 \end{split}
 \label{eq:continuation-image-loss}
\end{equation}
The sums cover every pixel of all $N$ complete targets, rather than balancing
each glyph independently. The final teacher-forced state $z_N$ predicts the
terminal boundary without an additional image target. The balanced stopping loss is
\begin{equation}
 \mathcal L_{\rm stop}=\frac12\left[
 \frac1N\sum_{t=0}^{N-1}\operatorname{softplus}(s_t)
 +\operatorname{softplus}(-s_N)\right].
 \label{eq:continuation-stop-loss}
\end{equation}
Training minimizes $\mathcal L_{\rm image}+0.1\mathcal L_{\rm stop}$.
Frozen encoder, decoder, source weights and affine statistics are not updated.
Image-loss gradients still pass through the frozen decoder into readout,
query and residual parameters. No decoded reference identity or lexical output
projection supplies these gradients.

At autonomous inference, evaluate stopping before another emission. If
$\sigma(s_t)\geq0.5$, end the answer; otherwise emit
\begin{equation}
 \widehat Y_{ti}=1-2\,\mathbf1\{\sigma(L_{ti})\geq0.5\}.
 \label{eq:continuation-emission}
\end{equation}
Only $S(E(\widehat Y_t))$ supplies the next input. An initially stopped stream
is empty; a stream reaching its image cap is retained as capped. Reference
length, font rendering, a tokenizer, OCR, lexical logits, a live teacher and an
answer bank are absent from this path. The physical pixel factorization and
teacher-forced objective do not by themselves guarantee semantic learning or
stability under generated histories.

\subsection*{Preflight evidence and the registered complete fit}
The real preflight selects four complete eligible training cases and excludes
the earlier feedback-ceiling diagnostic records and source groups. Its authentic
targets have 1, 74, 184 and 513 glyphs, totaling 772 positions. All 36 answer
layers receive actual gradients. An independent audit opens the saved checkpoint
and verifies four Adam steps for all 85 trainable tensors, FP32 moments,
original targets and prompt pixels. Prompt key/value, causality and frozen-body
version checks come from the registered GPU execution. Initial residual-down
tensors are not archived; their gradients, moments and bound update checks
are retained. Numerical split/native and serial/parallel comparisons use fixed
absolute tolerance 0.002 rather than asserting shape-independent equality.

Trainable parameters comprise 2,359,296 residual coordinates, a 4,096-coordinate
query and 17,848,833 readout/stopping parameters. The deployed system also retains
the 7,522,463,984-parameter inherited visual/contextual body and frozen image
interfaces. The 150.811-second preflight includes 20.020 seconds pressure waiting
and peaks at 18.122546\,GiB allocated memory; loading takes 113.573 seconds and
may overlap waiting. These measurements cover the four complete cases, not
arbitrary page sizes, complete foundation training or matched-quality efficiency.

The subsequent full trial starts fresh, without the preflight updates. Complete
eligible QA and book targets provide 2,036,247 positions per epoch. For a book
record of $n$ glyphs, expose the first
$\min(128,\max(1,\lfloor n/4\rfloor))$ glyphs and predict
the remainder. The 13,766 records are presented once per epoch for three epochs;
each gets one fixed-cell and two proportional prompts. Forty-five font/layout
rejections retain the complete original raster. All source pixels, target
boundaries, ordering, exclusions and evaluation controls are registered before
fitting. With one complete sequence per update, FP32 AdamW uses no weight decay,
peak learning rate $10^{-4}$, 500-update linear warm-up, cosine decay to $10^{-5}$
and global norm clipping at one. Initial and 10,324/20,649/30,973/41,298-update
snapshots are fixed. The 732 prescribed autonomous streams, output-first review,
complete-language pilot, retained successes and untouched confirmation determine
any subsequent language claim. The final proportional pilot requires at least
12/32 exact answers, 4/16 both-correct pairs, 24/32 readable answers and 4/12
coherent complete long-book continuations, retaining all eight union proportional
and both fixed exact cases. Each QA-erasure arm may score at most two exact
answers, with normal minus the better erased arm at least eight. At this update,
those quality measurements are
unmeasured; neither the preflight nor a finite training loss substitutes for them.

The single-device fit requires 48\,GiB available host memory and 20\,GiB free
GPU memory on admission, with 24\,GiB host reserve and a 22\,GiB allocated-device
ceiling at runtime. These host requirements are separate from the 24-GiB GPU.
The cumulative fit wall budget is 48 hours, including the failed first child's
139.558657 seconds. Pressure waiting is bounded within that budget. Offline data
preparation, inherited pretraining, earlier interface fits, evaluation and review
remain additional costs; feasibility does not establish an efficiency advantage.
